\section{Experimental Method}
\label{sec:evaluation}
The study asks three questions. How much of a burst reduction comes from work
placement within a fixed implementation? How does that choice interact with
native kernels and the completion horizon? And does it improve
complete-module behavior under Rack's sample-synchronous engine? The questions
and the workload matrix were recorded before collection. The analysis is a
\emph{descriptive pilot study}: it reports what was observed and does not
select a winner and then confirm it.

\subsection{Host, sessions, and workload boundaries}
Two separately prepared sessions ran on September 30, 2026 (local time) on one
Apple M1 Pro laptop with macOS 26.6.2 and Apple Clang 21.0.0, using the same
executable and archived sources. The operator reported disconnecting
networking and Bluetooth and closing agents and unnecessary applications. The
runner required AC power with Low Power Mode off, held \code{caffeinate -is}
assertions, settled for 180 seconds after staging, and settled each fresh
process for one second before warmup. Snapshots taken before and after each
session confirm the power and sleep assertions. They also show residual macOS
service activity, including \code{sandboxd} and Spotlight indexing, and cannot
rule out interference during any particular interval.

Each session ran 194 planned cells twice, each time in a fresh process, for
776 processes in total. Three bridge cells repeat across groups, leaving 191
unique configurations. Group order was fixed, and process order within each
group followed predeclared session seeds. Every process and every slow
interval is retained. Table~\ref{tab:study-design} summarizes the seven
groups. Setup, plan creation, correctness replay, storage probes, and output
serialization all fall outside the timed block intervals.

\input{.build/paper-study/design.tex}

The stream harness times either complete analysis calls feeding a controlled
sink or the modules' actual \code{process()} paths, including input
conditioning and display-data conversion. Fourier's shipped defaults use four
SIMD lanes, $N=2048$, a 30\,ms hop ($H=1440$ at 48\,kHz), and a flattop
window; Spectre's use scalar samples, $N=2048$, $H=1024$, and a flattop
window. Unless stated otherwise, controlled scalar comparisons use a Hann
window, $H=1024$, no smoothing, and $D=64$. Every table and figure states
which boundary it measures.

The host group drives the pinned Rack engine's actual \code{stepBlock()},
including worker synchronization, with one or four engine threads; 0, 1, 4,
or 16 analyzers; and 0, 16, or 64 deterministic background modules. These form
a structured set of combinations rather than a full factorial design. Native
module controls replace only the analysis inside the same module shell,
keeping input conditioning, display conversion, and mailbox publication. They
also retain some unused production storage, so their footprint does not
represent the memory requirement of an optimized replacement. An optional
concurrent headless consumer copies and checksums snapshots at 30 or 60\,Hz,
with declared stalls; it neither draws a window nor drives an audio interface.

\subsection{Matched controls and native layouts}
The matched first-party batch and distributed controls share one scalar
analysis implementation, its arithmetic, and its storage, so they isolate
placement more cleanly than any comparison across FFT libraries. Each native
batch/hybrid pair likewise shares one retained-input, window, magnitude,
smoothing, and output path. Batch runs the whole sequence immediately; hybrid
distributes the surrounding stages but executes the transform and its
reordering as one indivisible native call. The native pipelines read
contiguous ring segments and use native output layouts, avoiding per-element
ring-modulo arithmetic and conversion to a canonical complex layout. vDSP
uses vectorized windowing and packing and batched four-channel transforms.
FFTW uses reusable real plans, including \code{plan\_many} for four channels,
created in each fresh process by single-threaded \code{FFTW\_MEASURE} planning
without imported wisdom. PFFFT uses its ordered packed spectrum. All paths
compute magnitudes with an explicitly scaled norm to avoid premature underflow
when squaring. These are strong native controls, but they do not exhaust every
provider-specific optimization.

The granularity group holds the frame hop fixed and varies the native
completion horizon over $H_c\in\{H,\ceil{H/2},\ceil{H/4}\}$. Publication
occurs at endpoint age $H_c-1$, after which the pipeline only captures input
until the next endpoint. This trades latency against the density of
surrounding-stage work without subdividing the native FFT. Native sub-FFT
leaves and new stage weights were deliberately deferred and have no measured
evidence here. The production weights, $p=4$ for window rebuilds and $o=2$
for Fourier output, were tuned in earlier engineering experiments on the same
M1 Pro, and band-cache rebuilding runs inside output units without additional
weight. The study therefore does not independently validate those weights.

\subsection{Timing, replication, and age}
Stream groups execute blocks back to back. The host group releases blocks on
an absolute periodic schedule with period $D/f_s$ and does not reset the
release origin after an overrun. A monotonic clock records wake, start, and
finish times separately. Threads inherit the default scheduling policy; no
core affinity, frequency lock, or real-time priority was applied. Stream
processes inherit and record their floating-point control state, while the
actual Rack engine applies its own per-block FPU policy. Empty-timer
observations are retained without subtraction, and reported values are
rounded to about three significant figures.

The primary burst statistic is the median \emph{per-hop peak block time}: for
each complete interval from one frame endpoint to the next, take the longest
duration of any whole block that intersects it. Unlike a callback percentile,
this statistic always captures a batch burst, however rare such bursts are
among all callbacks. A block that spans two hops counts toward both, and no
sub-block costs are imputed. Engine peaks use actual replay endpoints and the
full hop for both immediate and delayed publication. Lifecycle sequences are
excluded from this stationary statistic. Callback p99, maxima, and counts of
blocks exceeding $D/f_s$ are reported as secondary observations.

Stream cost is the sum of instrumented block wall times divided by the number
of engine samples. It therefore includes timer, dispatch, and interruption
effects; the design has no separate untimed throughput pass. Host results also
record process-wide user-plus-system CPU time over the measurement interval,
including workers, pacing overhead, and any headless consumer. CPU time and
block wall time are distinct quantities. A \emph{release miss} means that a
block finished after its synthetic release deadline, including wake lateness
and accumulated backlog. A \emph{compute-budget exceedance} means that the
timed block alone exceeded $D/f_s$. Neither is an audio-device underrun.

For each cell, costs are averaged over the two processes within a session,
and burst quantiles take the median of the two process summaries. Tables
report the midpoint of the two session summaries, their range where
indicated, and the largest retained observation. There are only two session
replicates, both on the same date and host, and callbacks, hops, analyzers,
and channels are dependent observations. We therefore report descriptive
variation rather than confidence intervals, rare-event probabilities, or
worst-case execution-time claims~\cite{kalibera2013,mytkowicz2009,wilhelm2008}.

Publication age is measured in logical engine samples from the frame
endpoint; center age adds $(N-1)/2$. Consumer age is bracketed by the
completed-block watermarks observed around each copy, leaving one block of
uncertainty. Final drains are excluded from ordinary consumer ages. These
sample-clock ages are not wall-clock freshness during an overrun, and they
exclude repaint.

\subsection{Numerical and provenance checks}
Every session passed checks of archive and artifact hashes, the exact
declared matrix and process order, execution sidecars, and a numerical replay
of every output. The independent references are direct DFTs at small lengths
and a separately structured binary64 FFT at larger lengths; module replays
also model DC conditioning and display conversion. Relative $\ell_2$ and
$\ell_\infty$ errors of each magnitude vector must not exceed
$3\times10^{-4}$ in float or $10^{-10}$ in double, and a zero reference
requires zero output. These are engineering acceptance budgets, not
guarantees for every weak bin.

A long complete-module decay fixture separately flags the region affected by
Rack's flush-to-zero behavior and Fourier's display mapping. It applies the
predeclared \code{module-decay-ftz-v1} absolute floor, retains every vector,
and excludes flagged vectors from relative-accuracy claims. The floor is an
engineering allowance, not a proven error bound. The distinction matters
because relative error can reach one in that tail while absolute error stays
below the floor. Native allocators are opaque to the separate C++ allocation
probe, so their storage and allocation calls are not observed.
Appendix~\ref{app:reproducibility} identifies the raw evidence and the
deterministic publication tools.
